% !PRINTED 119
\weberpage{119}
Durch wiederholte Anwendung von (9) ergibt sich, wenn $a$, $b$ ganze Zahlen sind [vgl. \S\,17, (16)]:

\weq{11}{\sigma(u+a\omega_1+b\omega_2) = (-1)^{a+b+ab}\,e^{(a\eta_1+b\eta_2)(2u+a\omega_1+b\omega_2)}\,\sigma(u).}
Wenn man die Funktion $\sigma(u)$, wie sie durch die Formel (6) gegeben ist, nach Potenzen von $u$ entwickelt, so findet sich, daß nicht nur die zweite, sondern auch die dritte Potenz von $u$ in der Entwickelung nicht vorkommt, und man hat also:

\weq{12}{\sigma(0)=0,\quad \sigma'(0)=1,\quad \sigma''(0)=0,\quad \sigma'''(0)=0.}
Wenn auf $\omega_1, \omega_2$ eine lineare Substitution

\weq{13}{(\omega_1', \omega_2') = \begin{pmatrix}\alpha, & \beta\\ \gamma, & \delta\end{pmatrix}(\omega_1, \omega_2)}
angewendet wird, so erfahren die Größen $\eta_1, \eta_2$ die entsprechende Substitution

\weq{14}{(\eta_1', \eta_2') = \begin{pmatrix}\alpha, & \beta\\ \gamma, & \delta\end{pmatrix}(\eta_1, \eta_2),}
wie man unmittelbar aus den Relationen (7) und (11) folgert.

Differentiiert man die Formeln (9) logarithmisch nach $u$, so folgt

\weq{15}{\begin{aligned}
\frac{\sigma'(u+\omega_1)}{\sigma(u+\omega_1)} &= \frac{\sigma'(u)}{\sigma(u)} + 2\eta_1,\\
\frac{\sigma'(u+\omega_2)}{\sigma(u+\omega_2)} &= \frac{\sigma'(u)}{\sigma(u)} + 2\eta_2,
\end{aligned}}
und indem man in der ersten $u=-\dfrac{\omega_1}{2}$, in der zweiten $u=-\dfrac{\omega_2}{2}$ setzt, und beachtet, daß $\sigma(u)$ eine ungerade und folglich $\sigma'(u)$ eine gerade Funktion von $u$ ist:

\weq{16}{\eta_1 = \frac{\sigma'\left(\dfrac{\omega_1}{2}\right)}{\sigma\left(\dfrac{\omega_1}{2}\right)}.\quad \eta_2 = \frac{\sigma'\left(\dfrac{\omega_2}{2}\right)}{\sigma\left(\dfrac{\omega_2}{2}\right)}.}

\begin{center}
\textbf{\S\,36. \gesperrt{Die Funktionen $\sigma_{00}, \sigma_{01}, \sigma_{10}$.}}
\end{center}

Es bleibt uns noch übrig, die analogen Resultate für die übrigen Charakteristiken zu gewinnen. Wir gehen aus von der folgenden Bemerkung: Sind $x_1, x_2$; $y_1, y_2$ irgend zwei Paare von Größen, welche durch dieselbe lineare Substitution $S$ in $x_1', x_2'$; $y_1', y_2'$ transformiert werden, ist also:
%CONT>>
